% !TEX program = xelatex
% AI-effects 프로젝트 --- 한국 직업별 AI 활용도 2개(대화량 기반)의 계산 결과.
% 표 본문: results/table/25_*.tex (code/04_analysis_kor_ai_usage_index.do가 생성)
% 계산: code/01_import_kostat_regional_employment.do, 01_import_moel_wage_structure_hours.do,
%       03_merge_kor_conv_ai_ksco8.do(SOC 단계 절감 시간), 03_merge_kor_ksco3_ai_usage.do
%       → data/proc/merged/kor_ai_usage_ksco{3,2,1}
% 관련: report 25(자료와 준비사항), paper 19(한국 대화량), paper 21(연계표), paper 23(KSCO 결합)
% 컴파일: xelatex (kotex 필요)
\documentclass[11pt]{article}

\usepackage{fontspec}
\usepackage{kotex}
\usepackage[a4paper,margin=2.2cm]{geometry}
\usepackage{longtable}
\usepackage{booktabs}
\usepackage{array}
\usepackage{amsmath}
\usepackage{xcolor}
\usepackage{hyperref}
\usepackage{enumitem}

\hypersetup{colorlinks=true,linkcolor=blue!50!black,urlcolor=blue!50!black}
\newcolumntype{L}[1]{>{\raggedright\arraybackslash}p{#1}}
\newcolumntype{R}[1]{>{\raggedleft\arraybackslash}p{#1}}
\setlength{\LTleft}{0pt}
\setlength{\LTright}{0pt}
\renewcommand{\arraystretch}{1.15}
\setlength{\tabcolsep}{4pt}

\title{한국 직업별 AI 활용도\\[2mm]
\large --- 업무용 대화량과 지역별고용조사로 만든 활용도 2개의 계산 결과 ---}
\author{AI-effects 프로젝트}
\date{\today}

\begin{document}
\maketitle

\begin{abstract}
\noindent 한국 직업별 월간 업무용 AI 대화량(2026년 4$\cdot$5월, paper 19$\cdot$23)과 지역별고용조사 2022년 상반기 자료로, 사전적 지표인 AI 노출도와 구분되는 실현 지표 두 가지를 KSCO 8차 소분류$\cdot$중분류$\cdot$대분류 단위로 계산했다. 활용도 1은 취업자 1인당 월 업무용 대화 수, 활용도 2는 월 주업 근로시간 중 AI로 줄어든 시간의 비중이다. 대화당 메시지 수는 $K = 3$으로 두었고, 활용도 2의 분자(대화 수 $\times$ 대화당 절감 시간)는 SOC 단계에서 계산한 뒤 배분했다. 주요 결과는 다음과 같다.
\begin{itemize}[leftmargin=*]
  \item \textbf{덮개}: 소분류 163개(군인 제외) 중 129개에 지표가 있고, 이들이 2022년 취업자의 87.4\%를 차지한다. 대화량이 있는 직업은 모두 절감 시간이 있다.
  \item \textbf{수준}: 전국 활용도 1은 취업자 1인당 월 12.3회, 활용도 2는 31.6\%다. 활용도 2가 100\%를 넘는 소분류가 38개다. 대화 과업의 직업이 사용자의 직업과 다르고 절감 시간이 Claude의 추정치이기 때문이며, 활용도 2는 실제 절감 비율이 아니라 \textbf{순위 비교용}으로 써야 한다.
  \item \textbf{분포}: 관리자$\cdot$전문가에서 높고 농림어업$\cdot$단순노무에서 낮다. 소분류 상위는 생명$\cdot$자연과학 시험원, 통계$\cdot$데이터 사무원, 작가$\cdot$언론, 정보 보안, 정보 시스템 운영 등이다.
  \item \textbf{노출도와의 관계}: 노출도 지표 3개와의 순위상관은 0.58$\sim$0.66으로, 방향은 같지만 같은 지표는 아니다.
  \item \textbf{두 지표의 관계}: 활용도 1과 2의 순위상관은 0.97이다. log 활용도 2의 분산 중 82.5\%가 활용도 1에서, 17.9\%가 대화당 절감 시간에서 나온다.
  \item \textbf{강건성}: 분모 근로시간을 사업체 기준으로 보정하거나 부업 시간을 넣어도, 분모를 2025년으로 바꾸거나 달을 나눠도 순위상관은 0.995 이상이다.
\end{itemize}
\end{abstract}

\tableofcontents

% ============================================================
\section{목적과 정의}
% ============================================================

AI 노출도는 ``이 직업의 과업을 AI가 할 수 있는가''를 나타내는 사전적 지표다. 본 노트는 실제 사용량에서 출발한 직업별 \textbf{AI 활용도}를 만들고 그 결과를 정리한다. 자료와 절차의 설계는 report 25를 따랐다.

KSCO 8차 분류 단위 $u$(소분류$\cdot$중분류$\cdot$대분류)에 대해
\begin{align}
\text{활용도 1}_u &= \frac{C_u}{E_u}, &
\text{활용도 2}_u &= \frac{S_u}{H_u}\times 100,\qquad
S_u = \sum_s w^{\text{alloc}}_{s,u}\, C_s\, \frac{\tau_s}{60}
\end{align}
이다. $C$는 월 업무용 대화 수($K = 3$), $\tau_s$는 SOC $s$의 대화당 절감 시간(분, AEI 전 세계 값), $w^{\text{alloc}}$은 SOC $\to$ KSCO 배분 가중치(paper 21), $E$는 2022년 상반기 취업자 수, $H$는 2022년 상반기 월 총 주업 근로시간이다. 두 지표 사이에는
\[
\text{활용도 2}_u = \text{활용도 1}_u \times \frac{\tau^{\text{eff}}_u / 60}{4.345\, h_u}\times 100
\]
의 관계가 있다. $\tau^{\text{eff}}_u = 60\, S_u / C_u$는 대화량으로 가중한 대화당 실효 절감 시간, $h_u$는 1인당 주당 근로시간이다.

사용자 결정(2026-10-17)은 다음과 같다: 분모 = 지역별고용조사 2022년 상반기, 분석 단위 = 소분류(중분류$\cdot$대분류도 계산), $K = 3$, 활용도 2 분모 근로시간 = 지역별고용조사(다른 조사는 강건성 검토), 활용도 2 분자 = SOC 단계 계산.

% ============================================================
\section{자료와 계산}
% ============================================================

\subsection{분자: 업무용 대화량과 절감 시간}

\begin{itemize}[leftmargin=*]
  \item 대화량은 \texttt{kor\_ksco3\_conv\_ai\_indicators}(월 $\times$ 소분류)의 \texttt{conv\_k3}이다. 한국 8개 AI 제품의 업무용 메시지를 대화당 메시지 3개로 나눈 값을 AEI Claude.ai 한국 직업 구성으로 나누었다(paper 19).
  \item 절감 시간 곱은 \texttt{03\_merge\_kor\_conv\_ai\_ksco8.do}에서 SOC마다 $C_s \times \tau_s / 60$(\texttt{saved\_hr\_k3})을 만든 뒤 다른 양 변수와 같이 $w^{\text{alloc}}$으로 배분했다. KSCO 파일의 평균 절감 시간(\texttt{time\_saved\_min}, $w^{\text{mean}}$ 평균)과 대화량을 곱하는 방법은 쓰지 않았다(paper 23).
  \item \textbf{절감 시간 결측은 문제가 되지 않는다.} 절감 시간이 없는 SOC 195개는 AEI 전 세계 세부 직업 자료에 공개되지 않은 직업이고, 모두 한국 대화량도 공개되지 않았다. 한국 대화량이 공개된 SOC는 4월의 자리표시 행(\texttt{45-XXXX})을 빼면 모두 절감 시간이 있다(do 파일에서 확인). 그래서 배분되는 대화의 절감 시간 덮개는 100\%다.
  \item 두 달 처리: 소분류는 대화량이 있는 달만 평균했다. 중분류$\cdot$대분류는 달마다 소분류 분자를 더한 뒤(대화량 결측 소분류는 0) 평균했다. 미공개 직업의 대화는 이미 공개 직업에 비례 배분되어 있으므로 상위 분류의 합은 전국 합과 같다.
\end{itemize}

\subsection{분모: 지역별고용조사 2022년 상반기}

\begin{itemize}[leftmargin=*]
  \item MDIS 공공용 \textbf{지역별고용조사 상반기(B형, 전국, 소분류)} 2022년 자료(조사연월 2022년 4월). 2022년 자료도 직업이 KSCO 8차 소분류로 제공되어 7$\to$8차 변환이 필요 없다.
  \item 취업자(경제활동인구상태 1)만 쓰고 가중치는 시도전국가중값이다. 2022년 상반기 취업자는 2,808만 명, 주당 평균 주업 근로시간은 38.8시간이다. 소분류는 163개(군인 4개는 이 자료에 없음)다.
  \item 근로시간은 조사대상주간의 \textbf{주업} 근로시간이다. 부업은 어느 직업인지 알 수 없어 뺐다. 일시휴직자(근로시간 빈칸, 표본 2,572명)는 0시간으로 넣었다. 주당 시간은 $\times 4.345$로 월로 바꿨다.
  \item 매칭 키: \texttt{ksco3}(문자 3자리) = 지역별고용조사 \texttt{3자리\_8차직업소분류코드} = 대화량 파일 \texttt{ksco3}. 대화량 파일에만 있는 소분류는 군인 4개뿐이다.
\end{itemize}

\subsection{강건성 검토용 분모}

\begin{itemize}[leftmargin=*]
  \item \textbf{사업체 기준 근로시간 보정}: 고용형태별근로실태조사 2022(6월 기준) 임금근로자 월 실근로시간(소정 + 초과)을 지역별고용조사 임금근로자 월 주업시간으로 나눈 비율 $r$을 직업 대분류마다 구해 $H$에 곱했다. 공공용 자료에 직업이 대분류만 있어 대분류 안 소분류에는 같은 $r$를 쓴다. 장인성(2018) 3장처럼 비임금근로자에게도 같은 $r$를 적용했다. $r$는 기능원 0.80$\sim$단순노무 1.01이고 대부분 1보다 작다(가구조사 시간이 사업체 시간보다 김). 특수형태근로종사자는 2023년부터 직업$\cdot$시간이 비어 있어 모든 해에서 뺐다.
  \item \textbf{주업 + 부업 시간}, \textbf{2025년 상반기 분모}, \textbf{달별 값}(2026-04, 2026-05).
\end{itemize}

% ============================================================
\section{결과}
% ============================================================

\subsection{덮개와 전국 수준}

\begin{longtable}{L{1.6cm}R{1.3cm}R{1.6cm}R{1.8cm}R{1.7cm}R{1.8cm}R{1.6cm}R{1.6cm}R{1.6cm}}
\caption{분류 수준별 요약(전국값 = 분자 합 / 분모 합)}\label{tab:summary}\\
\toprule
 & 단위 수 & 지표가 있는 단위 & 취업자 덮개(\%) & 전국 활용도 1 & 전국 활용도 2(\%) & 중앙값 활용도 1 & 중앙값 활용도 2(\%) & 활용도 2 $>$ 100\% \\
\midrule
% 04_analysis_kor_ai_usage_index.do가 생성. 직접 고치지 말 것.
소분류 & 163 & 129 & 87.4 & 12.33 & 31.64 & 6.81 & 12.86 & 38 \\
중분류 & 56 & 50 & 95.7 & 12.32 & 31.64 & 3.82 & 4.96 & 10 \\
대분류 & 9 & 9 & 100.0 & 12.32 & 31.64 & 2.77 & 3.99 & 2 \\

\bottomrule
\end{longtable}

\begin{itemize}[leftmargin=*]
  \item 소분류 34개는 연결된 SOC의 대화량이 모두 미공개라 지표가 없다. 이들은 2022년 취업자의 12.6\%를 차지한다. 중분류와 대분류로 올라가면 덮개가 95.7\%, 100\%가 된다.
  \item 전국 활용도 1은 취업자 1인당 월 12.3회다(2026년 4$\cdot$5월 업무용 대화 수 $\div$ 2022년 취업자).
  \item \textbf{전국 활용도 2는 31.6\%이고, 소분류 38개는 100\%를 넘는다.} 절감 시간의 합이 그 직업 종사자의 근로시간보다 크다는 뜻이므로 실제 절감 비율로 읽을 수 없다. 원인은 \S\ref{sec:level}에서 정리한다.
\end{itemize}

\subsection{대분류}

\begin{longtable}{L{4.2cm}R{1.3cm}R{1.4cm}R{1.3cm}R{1.4cm}R{1.4cm}R{1.6cm}R{1.4cm}R{1.3cm}}
\caption{대분류별 활용도와 구성요소}\label{tab:ksco1}\\
\toprule
대분류 & 취업자(만 명) & 대화량 있는 소분류 & 취업자 덮개(\%) & 활용도 1 & 활용도 2(\%) & 활용도 2, 사업체 보정(\%) & 실효 절감 시간(분) & 주당 시간 \\
\midrule
% 04_06_analysis_kor_ai_usage_index.do가 생성. 직접 고치지 말 것.
1 관리자 & 42 & 16/16 & 100.0 & 4.57 & 5.50 & 5.56 & 336 & 41.7 \\
2 전문가 및 관련 종사자 & 585 & 45/48 & 99.1 & 3.38 & 3.50 & 3.41 & 266 & 38.4 \\
3 사무 종사자 & 485 & 13/13 & 100.0 & 0.85 & 0.82 & 0.80 & 255 & 39.4 \\
4 서비스 종사자 & 319 & 13/15 & 88.2 & 0.10 & 0.05 & 0.05 & 122 & 39.0 \\
5 판매 종사자 & 266 & 5/5 & 100.0 & 0.33 & 0.13 & 0.13 & 103 & 40.6 \\
6 농림∙어업 숙련 종사자 & 149 & 4/5 & 97.3 & 0.01 & 0.00 & 0.00 & 37 & 35.1 \\
7 기능원 및 관련 기능 종사자 & 242 & 14/19 & 70.5 & 0.22 & 0.15 & 0.18 & 191 & 41.3 \\
8 장치∙기계 조작 및 조립 종사자 & 307 & 10/30 & 52.8 & 0.08 & 0.08 & 0.08 & 286 & 42.8 \\
9 단순노무 종사자 & 412 & 9/12 & 78.2 & 0.06 & 0.03 & 0.03 & 105 & 33.7 \\

\bottomrule
\end{longtable}

\begin{itemize}[leftmargin=*]
  \item 관리자(취업자 1인당 월 56.3회)와 전문가(41.7회)가 가장 높고, 사무(10.4회), 판매(4.1회), 기능원(2.8회) 순이다. 서비스(1.2회), 장치$\cdot$기계 조작(1.0회), 단순노무(0.7회), 농림어업(0.06회)은 낮다.
  \item 대화당 실효 절감 시간은 관리자(336분), 장치$\cdot$기계 조작(286분), 전문가(266분), 사무(255분)에서 길고 농림어업(37분)에서 짧다. 그래서 장치$\cdot$기계 조작과 기능원은 활용도 1보다 활용도 2의 순위가 상대적으로 높다.
  \item 장치$\cdot$기계 조작(52.8\%)과 기능원(70.5\%), 단순노무(78.2\%)는 대화량이 있는 소분류의 취업자 덮개가 낮다. 이들 대분류의 값은 대화량 미공개 소분류를 0으로 더한 것이라 하한에 가깝다.
\end{itemize}

\subsection{상위 소분류}

\begin{longtable}{L{6.4cm}R{1.5cm}R{1.8cm}R{1.5cm}R{1.6cm}R{1.8cm}}
\caption{활용도 1 상위 15개 소분류}\label{tab:top1}\\
\toprule
소분류 & 취업자(천 명) & 월 대화 수(만 건) & 활용도 1 & 활용도 2(\%) & 실효 절감 시간(분) \\
\midrule
% 04_analysis_kor_ai_usage_index.do가 생성. 직접 고치지 말 것.
213 생명 및 자연과학 관련 시험원 & 11 & 964 & 839.8 & 2720.0 & 337 \\
332 통계‧데이터 관련 사무원 & 16 & 904 & 569.7 & 2353.6 & 400 \\
149 기타 건설‧전기 및 생산 관련 관리자$^{b}$ & 3 & 146 & 567.3 & 2181.4 & 400 \\
291 작가 및 언론 관련 전문가 & 96 & 4,126 & 430.0 & 1028.5 & 228 \\
292 학예사‧사서 및 기록물 관리사 & 31 & 1,318 & 422.6 & 676.3 & 156 \\
122 마케팅 및 광고‧홍보 관리자$^{b}$ & 9 & 350 & 399.5 & 1465.5 & 385 \\
223 네트워크 및 정보 보안 전문가 & 33 & 1,328 & 397.8 & 934.9 & 253 \\
225 정보 시스템 및 웹 운영자 & 79 & 2,878 & 365.2 & 580.1 & 166 \\
135 정보 통신 관련 관리자 & 6 & 188 & 297.0 & 319.2 & 115 \\
298 문화‧예술 관련 기획자 & 10 & 229 & 234.8 & 863.4 & 385 \\
283 상품 기획‧홍보 및 조사 전문가 & 62 & 1,453 & 233.7 & 724.1 & 318 \\
159 기타 판매 및 고객 서비스 관리자$^{b}$ & 1 & 29 & 224.4 & 776.8 & 383 \\
153 환경‧청소 및 경비 관련 관리자$^{b}$ & 4 & 93 & 212.9 & 609.4 & 316 \\
212 인문 및 사회과학 전문가 & 22 & 438 & 200.8 & 470.1 & 230 \\
224 데이터 전문가 & 25 & 380 & 152.2 & 611.9 & 430 \\

\bottomrule
\multicolumn{6}{l}{\footnotesize $^{a}$ 소분류 안 세분류 일부만 대화량 있음(하한값). $^{b}$ 지역별고용조사 표본 30 미만.}
\end{longtable}

\begin{longtable}{L{6.4cm}R{1.5cm}R{1.8cm}R{1.5cm}R{1.6cm}R{1.8cm}}
\caption{활용도 2 상위 15개 소분류}\label{tab:top2}\\
\toprule
소분류 & 취업자(천 명) & 월 대화 수(만 건) & 활용도 1 & 활용도 2(\%) & 실효 절감 시간(분) \\
\midrule
% 04_analysis_kor_ai_usage_index.do가 생성. 직접 고치지 말 것.
213 생명 및 자연과학 관련 시험원 & 11 & 964 & 839.8 & 2720.0 & 337 \\
332 통계‧데이터 관련 사무원 & 16 & 904 & 569.7 & 2353.6 & 400 \\
149 기타 건설‧전기 및 생산 관련 관리자$^{b}$ & 3 & 146 & 567.3 & 2181.4 & 400 \\
122 마케팅 및 광고‧홍보 관리자$^{b}$ & 9 & 350 & 399.5 & 1465.5 & 385 \\
291 작가 및 언론 관련 전문가 & 96 & 4,126 & 430.0 & 1028.5 & 228 \\
223 네트워크 및 정보 보안 전문가 & 33 & 1,328 & 397.8 & 934.9 & 253 \\
298 문화‧예술 관련 기획자 & 10 & 229 & 234.8 & 863.4 & 385 \\
159 기타 판매 및 고객 서비스 관리자$^{b}$ & 1 & 29 & 224.4 & 776.8 & 383 \\
283 상품 기획‧홍보 및 조사 전문가 & 62 & 1,453 & 233.7 & 724.1 & 318 \\
292 학예사‧사서 및 기록물 관리사 & 31 & 1,318 & 422.6 & 676.3 & 156 \\
224 데이터 전문가 & 25 & 380 & 152.2 & 611.9 & 430 \\
153 환경‧청소 및 경비 관련 관리자$^{b}$ & 4 & 93 & 212.9 & 609.4 & 316 \\
225 정보 시스템 및 웹 운영자 & 79 & 2,878 & 365.2 & 580.1 & 166 \\
212 인문 및 사회과학 전문가 & 22 & 438 & 200.8 & 470.1 & 230 \\
269 기타 교육 전문가 & 46 & 489 & 106.9 & 453.5 & 372 \\

\bottomrule
\multicolumn{6}{l}{\footnotesize $^{a}$ 소분류 안 세분류 일부만 대화량 있음(하한값). $^{b}$ 지역별고용조사 표본 30 미만.}
\end{longtable}

\begin{itemize}[leftmargin=*]
  \item 두 목록의 구성은 거의 같다. 상위는 연구$\cdot$데이터(213 생명$\cdot$자연과학 시험원, 332 통계$\cdot$데이터 사무원, 224 데이터 전문가), 글쓰기$\cdot$기획(291 작가$\cdot$언론, 283 상품 기획$\cdot$홍보, 298 문화$\cdot$예술 기획자), 정보통신(223 정보 보안, 225 정보 시스템$\cdot$웹 운영자, 135 정보 통신 관리자)이다.
  \item \textbf{213 생명$\cdot$자연과학 시험원이 1위인 것은 연계 경로의 영향이 크다.} paper 23에서 보았듯이 세분류 2133(자연과학 시험원) 대화의 69\%는 ISCO 3314(통계$\cdot$수학 준전문가)를 거쳐 Data Scientists(15-2051) 등에서 들어온다. 취업자는 11천 명으로 작은데 대화는 월 964만 건이 배분되어 활용도가 커졌다.
  \item 관리자 소분류 몇 개(149, 122, 159, 153)는 표본이 30 미만이다($^{b}$). 이 값은 표본오차가 크다.
\end{itemize}

\subsection{노출도와의 관계}

\begin{longtable}{L{3.8cm}R{1.8cm}R{1.8cm}R{2.2cm}R{2.2cm}R{2.2cm}}
\caption{소분류 순위상관(Spearman, 지표가 모두 있는 소분류 129개)}\label{tab:corr}\\
\toprule
 & 활용도 1 & 활용도 2 & Observed exposure & Felten AIOE(LM) & Eloundou GPT-4 $\beta$ \\
\midrule
% 04_06_analysis_kor_ai_usage_index.do가 생성. 직접 고치지 말 것.
상대 활용도 1(RU1) & 1 &  &  &  &  \\
상대 활용도 2(RU2) & 0.98 & 1 &  &  &  \\
Observed exposure & 0.59 & 0.59 & 1 &  &  \\
Felten AIOE(LM) & 0.59 & 0.61 & 0.78 & 1 &  \\
Eloundou GPT-4 $\beta$ & 0.66 & 0.66 & 0.83 & 0.88 & 1 \\
% N(모든 지표가 있는 소분류) = 129

\bottomrule
\end{longtable}

\begin{itemize}[leftmargin=*]
  \item 활용도와 노출도의 순위상관은 0.58$\sim$0.66이다. 노출도가 높은 직업에서 실제 사용도 많지만, 노출도 지표끼리의 상관(0.78$\sim$0.88)보다 낮다. 활용도는 노출도와 다른 정보를 담고 있다.
  \item Observed exposure는 같은 AEI 자료에서 나왔는데도 상관이 가장 높지는 않다(0.58$\cdot$0.59). Eloundou GPT-4 $\beta$와의 상관이 가장 높다(0.65$\cdot$0.66).
  \item 노출도 지표는 KSCO에서 SOC를 $w^{\text{mean}}$으로 평균한 값이고 활용도는 $w^{\text{alloc}}$으로 배분한 대화량에서 나왔으므로, 같은 SOC 정보라도 가중 방식이 다르다.
\end{itemize}

\subsection{활용도 2의 분해}

\begin{longtable}{L{7.2cm}R{2.2cm}R{3.0cm}}
\caption{log 활용도 2의 분산 분해(소분류 129개; 기여 = 공분산(구성요소, log 활용도 2) / 분산(log 활용도 2))}\label{tab:decomp}\\
\toprule
구성요소 & 표준편차 & 분산 기여(\%) \\
\midrule
% 04_analysis_kor_ai_usage_index.do가 생성. 직접 고치지 말 것.
log 활용도 1(취업자 1인당 대화 수) & 2.48 & 82.5 \\
log 대화당 실효 절감 시간 & 0.74 & 17.9 \\
$-$ log 1인당 주당 근로시간 & 0.13 & -0.3 \\
\midrule
log 활용도 2 & 2.94 & 100.0 \\
% N = 129

\bottomrule
\end{longtable}

\noindent 활용도 2의 직업 간 차이는 대부분(82.5\%) 취업자 1인당 대화 수에서 나온다. 대화당 실효 절감 시간이 17.9\%를 설명하고, 1인당 근로시간의 차이는 거의 기여하지 않는다($-0.3$\%). 그래서 두 지표의 순위상관이 0.97로 높다. 활용도 2를 따로 쓰는 의미는 절감 시간이 긴 과업(장치$\cdot$기계 조작, 관리자 등)의 비중을 반영한다는 데 있다.

\subsection{강건성}

\begin{longtable}{L{6.2cm}R{1.2cm}R{1.8cm}R{1.8cm}R{1.6cm}R{1.6cm}}
\caption{기준 대비 대안의 소분류 순위상관과 비(대안 / 기준)}\label{tab:robust}\\
\toprule
비교 & $N$ & 순위상관 & 비 중앙값 & 비 최소 & 비 최대 \\
\midrule
% 04_06_analysis_kor_ai_usage_index.do가 생성. 직접 고치지 말 것.
RU2: 분모 사업체 기준 보정 & 129 & 0.999 & 0.979 & 0.914 & 1.152 \\
RU2: 분모 주업+부업 시간 & 129 & 1.000 & 1.000 & 0.963 & 1.002 \\
RU1: 분모 2025년 상반기 & 129 & 0.995 & 0.972 & 0.510 & 2.175 \\
RU2: 분모 2025년 상반기 & 129 & 0.995 & 0.971 & 0.483 & 1.786 \\
RU1: 2026-04 대 2026-05 & 129 & 0.991 & 0.982 & 0.333 & 2.240 \\
RU2: 2026-04 대 2026-05 & 129 & 0.996 & 0.989 & 0.332 & 1.599 \\

\bottomrule
\end{longtable}

\begin{itemize}[leftmargin=*]
  \item \textbf{근로시간 측정}: 사업체 기준으로 보정하면 활용도 2가 중앙값 기준 6.5\% 커지지만(가구조사 시간이 김) 순위는 거의 그대로다(0.999). 부업 시간을 넣어도 마찬가지다.
  \item \textbf{기준 연도}: 분모를 2025년 상반기로 바꿔도 순위상관은 0.995다. 비의 범위(0.48$\sim$2.11)는 2022$\sim$2025년 사이 고용이 크게 변한 소분류가 있음을 보여 주지만 순위에는 영향이 작다. report 25의 권장대로 노동시장 효과 분석에는 2022년 기준(사전적 분모)을 쓴다.
  \item \textbf{달}: 4월과 5월의 순위상관은 0.996$\cdot$0.997이다. 소분류 122개가 두 달 모두 대화량이 있다.
\end{itemize}

% ============================================================
\section{활용도 2의 수준을 해석할 때 주의할 점}\label{sec:level}
% ============================================================

활용도 2가 100\%를 넘는 소분류가 38개이고 전국값이 31.6\%다. 실제로 근로시간의 3분의 1이 줄었다고 볼 근거는 없으므로, 수준이 부풀려진 원인을 정리한다(report 25 4.1절).
\begin{enumerate}[leftmargin=*]
  \item \textbf{대화 과업의 직업 $\neq$ 사용자의 직업}: AEI는 대화 내용이 어느 직업의 과업인지를 분류한다. 데이터 분석을 하는 사무직이나 학생의 대화도 ``통계 전문가'' 과업으로 잡힌다. 분자는 그 과업의 전체 사용량이고 분모는 그 직업 종사자만이므로, 과업이 여러 직업에 퍼져 있을수록 활용도가 부풀려진다.
  \item \textbf{절감 시간은 Claude의 추정치}: 대화당 ``사람 혼자 걸리는 시간 $-$ AI와 함께 걸리는 시간''을 Claude가 추정한 값의 평균이며 실측이 아니다. 범위도 전 세계 Claude.ai다. 대화 하나가 실제 업무 하나를 대신했다는 가정도 들어 있다.
  \item \textbf{대화량 구성 가정}: 한국 8개 제품의 업무용 대화를 Claude.ai 한국 직업 구성으로 나눴다(paper 19의 B1). Claude는 코딩$\cdot$분석 비중이 높아 해당 직업으로 대화가 쏠린다.
  \item \textbf{연계 경로}: SOC 대화가 ISCO를 거쳐 KSCO 소수의 직업으로 몰리면(예: 213) 그 직업의 활용도가 커진다(paper 23).
\end{enumerate}
\noindent 따라서 활용도 2(와 활용도 1)의 수준은 실제 효과의 크기가 아니라 직업 간 \textbf{상대 순위}를 보는 데 쓴다. $K$를 바꾸면 수준은 $1/K$배로 달라지지만 순위는 변하지 않는다.

% ============================================================
\section{산출 파일과 다음 단계}
% ============================================================

\begin{itemize}[leftmargin=*]
  \item \texttt{data/proc/merged/kor\_ai\_usage\_ksco3.\{dta,csv\}}: 소분류 163개. 활용도 1$\cdot$2, 강건성 변수, 구성요소(\texttt{tau\_eff}, \texttt{hrs\_wk}), 노출도 지표, 플래그(\texttt{conv\_lb}, \texttt{small}).
  \item \texttt{kor\_ai\_usage\_ksco2}, \texttt{kor\_ai\_usage\_ksco1}: 중분류 56개, 대분류 9개. 노출도 지표는 2022년 취업자 가중 평균.
  \item \texttt{results/table/25\_kor\_ai\_usage\_index.xlsx}: 세 수준의 전체 변수.
\end{itemize}

다음 단계 후보:
\begin{enumerate}[leftmargin=*]
  \item 활용도(2022년 기준)와 2022$\to$2025년 직업별 고용$\cdot$근로시간$\cdot$임금 변화의 관계(지역별고용조사 2023$\sim$2025년 자료 보유).
  \item 2026년 상반기 지역별고용조사가 공표되면 2026년 기준 활용도(기술 통계용)를 추가.
  \item 대화 과업 직업과 사용자 직업의 차이를 줄이는 보정(예: 한국 근로자 AI 이용 설문의 직업별 이용률과 비교).
  \item 산업별 LCE(report 24)로의 확장: 활용도를 산업 $\times$ 직업 고용 분포로 산업에 배분.
\end{enumerate}

\end{document}
